
\begin{obeassessment}
\textbf{Kuis Bab 15}
1. Jelaskan apa yang dimaksud dengan \textit{One-Way Function} dan berikan contoh penerapannya.
2. Bandingkan antara SHA-256 dan MD5 dari sisi panjang output dan tingkat keamanan.
3. Bagaimana mekanisme kerja HMAC dalam menjamin integritas sekaligus otentikasi pesan?
4. Tuliskan aturan padding Merkle--Damgård dan hitung berapa bit nol yang disisipkan untuk pesan sepanjang 1000 bit. Berapa blok yang terbentuk?
5. Mengapa ketahanan kolisi sebuah fungsi hash hanya sebesar $2^{n/2}$? Turunkan hubungan itu dari paradoks ulang tahun.
6. Sebutkan nilai awal (IV) SHA-256 dan jelaskan dari mana kedelapan nilai itu diperoleh.
7. Jelaskan mengapa skema $H(K \| m)$ tidak aman, sedangkan $H(m \| K)$ gagal karena alasan yang berbeda. Serangan apa yang menjatuhkan masing-masing?
8. Bandingkan HMAC, CMAC, dan Poly1305 dari sisi bahan dasar, kecepatan, dan lingkungan pemakaian yang cocok.
9. Apa yang dimaksud dengan AEAD? Mengapa ketiga skema AEAD yang dibahas menuntut nonce yang unik, dan apa akibat pelanggaran syarat itu pada GCM?
10. Jelaskan mengapa perbandingan tag yang berhenti pada byte pertama yang berbeda membahayakan, dan tuliskan bentuk pemeriksaan yang aman.
\end{obeassessment}
\begin{competencychecklist}
\begin{tabularx}{\linewidth}{|X|c|}
\hline
Indikator Kompetensi & Tercapai \\
\hline
Saya dapat membedakan fungsi hash dengan enkripsi & [ ] \\
\hline
Saya dapat menjelaskan tiga properti utama hash kriptografis & [ ] \\
\hline
Saya dapat menurunkan batas ulang tahun dan menerapkannya pada MD5 dan SHA-256 & [ ] \\
\hline
Saya dapat menjelaskan aturan padding Merkle--Damgård beserta akibatnya & [ ] \\
\hline
Saya dapat membandingkan algoritma MD5, SHA-1, SHA-2, dan SHA-3 & [ ] \\
\hline
Saya dapat menelusuri jejak kompresi SHA-256 dan memeriksanya dengan Python & [ ] \\
\hline
Saya dapat menjelaskan serangan panjang-ekstensi dan cara menutupnya & [ ] \\
\hline
Saya dapat menguraikan cara kerja MAC dan HMAC & [ ] \\
\hline
Saya dapat membandingkan HMAC, CMAC, GMAC, dan Poly1305 & [ ] \\
\hline
Saya dapat menjelaskan prinsip AEAD dan syarat keunikan nonce & [ ] \\
\hline
Saya dapat menganalisis risiko serangan kolisi pada fungsi hash & [ ] \\
\hline
Saya dapat menuliskan pemeriksaan tag yang tahan terhadap serangan waktu & [ ] \\
\hline
\end{tabularx}
\end{competencychecklist}
